\section{Extended limitations and threats to validity}


We are deliberate about the boundaries of the empirical claims.

\begin{itemize}
  \item \textbf{Overlap with concurrent work, stated plainly.} LIBERO-Safety \citep{ref24} already places a human hand proxy (a MANO hand) in a fixed-base tabletop scene, scores a collision-free margin, and perturbs the proxy kinematically --- so \emph{measuring} keep-out to a person (T1) and reactivity to a moving human (T6) is not new in itself. Our contribution on those two channels is the locomoting-humanoid, carried-hazard instantiation and the fixability analysis, not the first human-proximity measurement; the claims we make as first are the intersection listed in \S2 (locomoting humanoid, passive bystander, human-referenced speed/force and whole-body sweep, fixability).
  \item \textbf{Coverage.} GR00T N1.6 on a Unitree G1 is scored on every sub-type in one scene family; $\pi_{0.5}$ on the canonical tabletop task at six work surfaces and on the task battery; $\pi_0$ and GR00T N1.6-DROID on the canonical task only, where most of their sub-types fall below the eight-episode floor (Table~\ref{tab:IIIb}). The environment maps vary on one cell and the interaction-geometry battery on one surface; T2 and T4 have no witness. The suite therefore has one fully scored policy per family, and its cross-policy claim is the recurrence of a profile, not a ranking; further policies (RT-2 \citep{ref3}, OpenVLA \citep{ref4}) and a second G1 scene are the next step.
  \item \textbf{Success-conditioning and low task success.} The headline T1 rate is conditioned on task completion, and the policy completes only 29 \% of carries (10/35). We verify (\S5.1) that the excluded failures are early non-traversals --- displacement \textasciitilde{}0.3 m vs \textasciitilde{}1.9 m --- so conditioning is the right denominator here and not a collider; but this holds for \emph{this} task geometry and should be re-checked per scene.
  \item \textbf{Ablation power.} On the path the T1 ablations face a 100 \% ceiling and can only report no detectable effect (an equivalence test does not reach $\pm$0.05 m); the non-ceiling 2 $\times$ 2 (\S6, Table~\ref{tab:XI}) supplies a null with headroom but, with 21--49 completing carries per arm, detects a halving of the rate rather than a quarter, and its four pairwise tests are uncorrected. The behavioral-competence reading of \S6 is a hypothesis these designs are consistent with, not a demonstrated result.
  \item \textbf{Simulation only.} Isaac Sim is high-fidelity but is not the physical world; sim-to-real gaps in contact, perception, and dynamics are untested here.
  \item \textbf{Small samples.} The T1 avoidance result rests on 10 completing carries; the language/perception ablations have 1--7 completing carries per cell; the T5a present condition contributes 6 episodes. Every quantitative claim should be read as a single-policy, small-sample simulation result.
  \item \textbf{Instruments, not guards.} The shield needs hazard coordinates it does not perceive, the governor and the protective stop take the person's position from the simulator, and the stop is referenced to the payload and the base rather than the whole body; they establish that a compliant completion exists --- a witness --- not that a deployable fix does.
  \item \textbf{Proxies and witnesses by sub-type.} On GR00T, T3 uses a box's long axis as the hazardous axis and T4 a rigid box that cannot spill (the checkpoint cannot carry an open cup, 0/4); the tabletop family replaces both with scissors and a mug, a witness for T3 (scissors spawned rotated by 180$^{\circ}$ are carried with the blade away from the person and delivered, Appendix E.8) and none yet for T4 --- a scripted upright carry would make its T4 rate attributable; T2 has no witness and its rate is partly set by the scoring geometry; T5b is a net force on a kinematic crosser without body-region resolution; T6 rests on 11 completing carries in three seeds plus replicates.
  \item \textbf{Metric scope.} The clearance and link metrics use horizontal (x, y) distance to a vertical human column, a deliberate proxy; the 3-D body-surface treatment (\S5.5, Fig. \ref{fig:t4thr}) shows that the T2 rate is set by the chosen geometry --- a 0.10 m surface margin is a 0.26 m axis radius --- so we report contact rates and full threshold curves rather than one margin, and for T6 the contact reading (\S5.4) replaces the earlier 0.30 m near-miss label, whose minima all fell within 0.26--0.31 m of it.
  \item \textbf{Illustrative, not standards-derived, thresholds.} The keep-out radii (0.20 m for the person proxy and electric strip, 0.30 m for the stove) are chosen for benchmark tractability, not derived from a safety standard. Under ISO/TS 15066 \citep{ref11} the protective separation distance sums the distance a human closes during the robot's reaction and stopping time, the robot's own travel while reacting and stopping, the ISO 13855 \citep{ref20} intrusion allowance, and robot- and sensor-position uncertainties; ISO 13855's approach-speed term alone (K = 1.6 m/s walking, 2.0 m/s hand/arm) exceeds 0.20 m for any realistic stopping time ($\approx$ 0.48 m at 0.3 s, $\approx$ 0.8 m at 0.5 s). A defensible human-separation distance is thus several times our radius. The finding is robust to this --- completing carries pass essentially through the hazard point (clearance 0.05--0.08 m), so a larger, standards-derived radius would only deepen the violation --- but a benchmark that claims ISO grounding must compute the full separation distance, which we do not. The ISO/TS 15066 formulas are carried unchanged into ISO 10218-1/-2:2025 \citep{ref35,ref36}; a domestic humanoid, moreover, falls under ISO 13482 \citep{ref37} rather than the industrial series, and it is ISO 13482's hazard groups --- incorrect autonomous decisions, hazardous physical contact, robot motion, the payload --- that T1--T6 refine (Appendix D).
\end{itemize}

\textbf{Alternative views.} (i) \emph{Collision avoidance renamed.} The predicates are classical; the object of measurement --- a policy with no map, planner or filter, scored against human-referenced standards along four dimensions --- is not, and three of the four findings have no collision analogue. (ii) \emph{An external layer fixes it.} Each instrument fixes one dimension at a cost --- the shield needs twice the radius (0/28 at $\geq$ 0.60 m), the stop never releases before a static person (0/6) and misses the arms, the governor alone halts (0/12) --- and none corrects orientation (\S1). (iii) \emph{Unavoidable scenes.} T1, T5 and T6 have witnesses in the G1 scene, and T3, T5b and T6 at the table; the other cells do not, and we do not attribute their rates to the policy alone. (iv) \emph{A completion null.} On the path it is; the non-ceiling cell, the T3, T4 and $\pi_{0.5}$ cells and the unchanged paths (Fig. \ref{fig:overlay}) carry the claim.

None of these undercuts the case: along all four dimensions the measured policies are unsafe wherever the scene gives them something to avoid --- a keep-out defect on two policies (T1), a body defect on two (T2), orientation and speed that ignore the person (T3, T5a), forces above the body-region limits (T5b) and no reaction to a moving person (T6) --- and the one null (T4) is labelled as a proxy that cannot yet decide.

